\section{The Future Order of Nations}

\begin{quotationx}
In the September, 1941 issue of \textit{La Vita Italiana}, \textbf{Julius Evola} wrote a review of Das Reich der Völker by the German political theorist, \textbf{Hans Keller}. The first part of Evola's essay is a review of the book. In the second part, whose translation appears below (\emph{ff}), Evola provides his own conception of the role of the Volk. Clearly, Evola cannot support Keller's naturalism nor the role of the Volk as the vanguard or the base for the political organization.

In the Traditional perspective, the Volk is lunar or passive and therefore has to rely on a spiritual and political hierarchy for its full development. Thus Evola must needs oppose any sort of nationalism on the basis of an ethnicity or a race. In rejecting that sort of neo-Nordicism, Evola will lead us back to the ideal of Romanity.
\end{quotationx}

\paragraph{Part I} \emph{In primis et ante omnia} we cannot adhere to the mythology of the \emph{Volk}, of the ``people nation'', so dear to Keller, who makes them the base of his juridical edifice. Mussolini once spoke ironically of the ``mysterious entity that calls itself the people''. ``People'', be it only as ``Volk'', is a simple myth that always and inevitably sounds like demagogy: of the type which is accompanied by intense polemical applications to disvalue and degrade the significance of everything that is the State, the formative political force from above. Keller stands on the most naïve optimistic natural law: he believes that the people exist as a very precise entity, equipped with its own consciousness, its own will, determined by ``eternal laws'' superior to all political forms in which they are concretized, a depositary of determinate values. We can only speak of similar things to some naïve types for whom, today, the natural meaning of nationality (quite different from every ``nationalism'' and deprived of a ``political'' character) becames something extremely tenuous in the face of becoming the ``mass'' of the people and the advent of a civilization no longer based on truly traditional values. Hegel said, ``the people is that part of the State that does not know what it wills''. That is precise. Our fascist idea is that the people, the nation, exists only as the State, in the State, and in a certain measure, only through the work of the State.

But our State is not the end of Keller's artificial antitheses: it is not a juridical superstructure, a mere fact of ``power'', an external power without basis. Our State is ethical and spiritual. It has the value of an \emph{entelechy}, i.e., of a formative force on the nation and the ``people'', who otherwise would remain a diffuse and unformed reality, vegetating on a naturalistic plane of life, without any metaphysical, ethical, or truly heroic tension.

Our point of view is detached both from the idea of natural law and the collectivizing of the people, as well as the abstract, juridical and rationalistic of the State. We are realists. We do not believe in the ``people'', instead we believe in a guiding, formative elite and, wherever it happens, leaders of the people. Keller has in mind the State only as \emph{caput mortuum}, i.e., as that which became in some cases, when the political structure created by leaders and elites personifying living traditions depersonalized themselves, objectified themselves, created an excuse to justify themselves, giving to understand the existence of public autonomous powers and of neutral juridical forms or norms to order to obtain in that way recognition to which the direct affirmation and the prominent prestige by the high stature of the Leaders and by superior princes was not enough.

To reach the same goal, i.e., to hold it without exhibiting it, along other directions, therefore, it is joined to the creation of a myth opposite in appearance to that of the neutral ``State'', but in reality answering to the same end: the myth of the ``people''. Keller lets himself avoid the admission that the idea of the Volk was ``discovered'' in 1933. Even if the date is not for us so recent, having some precise antecedents especially in democratic nationalism and French Jacobinism, (in fact even in France, as the antitheses to an excess of \emph{statism}, determined by the centralizing, absolutistic, and anti-aristocratic work followed by the kings up to Philip the Fair, it left, for the first time, the mystique of the nation, of the people as source of every law, not putting up with any authority from above)---even, therefore, if the Volk was not even discovered in 1933, only it stands in fact that the mythology referring to it is new, it was unknown by preceding civilizations and there came to life for precise political and propagandistic reasons: the ``people'' and \emph{Volk}, far from being notions having, today, a real content, are two idea-forces, two myths taken up by a power to affirm a given political system, determined by circumstances, to capture and organize the forces---in themselves directionless and apt to follow very different suggestions---of the true people and to strengthen therefore a given type of political authority. For the less important socializing tendencies, it was not necessary to us in Italy to have recourse to this myth: to us the spiritual idea of the State, and moreover, the direct authority that proceeds from a Duce and from a Monarch were a sufficient base to lead to the same result.

Keller says that all the differences between \emph{statute} law and law based on the idea of Volk lie in the judgment on human nature: in the first case, as in every idea of Reich and Imperium of Roman Catholic origin, it is pessimistic and does not have faith in human nature, while in the second case it has. But what is this ``human nature''? We are not pessimists, but realists. We believe in human nature, but not of everyone, less than ever in the collectivity as the mass, whose psychology we know well. Instead, we believe in that of the minority who create the States, who animate the peoples, who lead the collectivities to the heights which they could never reach alone. Neither statism, therefore, nor natural law, but the aristocratic hierarchical ideal, without the charade that is necessary only to demagogues and weak natures. That is what we mean by Roman and Aryan realism.

Keller says: ``to the Romanesque people, only the State is important, even if they serve nationalistic motives in addition.''

We say that we can even renounce the additive motives; we value the Mussolinian idea, according to which the people is something weak and blind, before they came to constitute a reality and a unitary will, compenetrated by higher meanings with the birth of the State. But we are not statists, because we do not make an idol of the State, as Keller and certain German jurists make of the Volk: because behind the State, in our view, there is somebody, there are Leaders, Monarchs---if you want, there is a ``superrace'', in which alone ``nation'', ``people'', ``race'', ``tradition'' cease to be abstractions. Every great ``people'' is always comprised of various influences, various racial elements, various traditions. It is the work of the elites to choose, to affirm, from among them all, a given element, to subordinate every other to it, leading back in that way to a precise order which otherwise would have remained heterogeneous and arrested in the form of a confused potentiality.

\paragraph{Part II} 


\begin{quotationx}
As Evola continues his analysis of the German legal scholar Hans Keller, we see recognizable themes. In particular, Keller comes surprisingly close to the New Right idea of Ethnopluralism. In Klemmer's scheme, each people is free to define themselves, to realize themselves, and live independently of each in their own enclaves. So, the ``New Right'' is not so new after all, as it merely resurrects stillborn ideas from the past. Evola rejects this conception on several grounds.

The first and most obvious is Keller's naturalistic starting point. Evola points out that for every Aryan civilization, the authority of their law was based on its relationship to the supernatural, and the lawgivers themselves were considered ``divine''. So Klemmer, and by extension the New Right, have nothing at all to do with the self-understanding of the European peoples of the past. Keller regards such claims as fantasies and projections of earthly leaders onto an imaginary plane. Any rightist movement must be related to a transcendental tradition; the pseudo-tradition of neo-paganism would not qualify for reasons that will be made clear in Part III.

The second difficulty with Keller's thesis is the very problem of defining a ``volk''. Keller speaks of nations of destiny or becoming, without specifying their actual destiny or the goal of the becoming. As such, it is not very helpful. Furthermore, it implies that there will be nations in a lesser or greater state of becoming. Would the latter have to treat the former as minors and guide their development?

In common with similar ideas held by new right movements today, Keller blamed ``inorganicity and leveling'' on the preceding religious Regnum. This thesis, Evola rejects. Keller's ethnopluralism is itself deracinated, since it lacks any content, or principles, that a supernatural conception of the Regnum would provide. Relations between ethnocentric conclaves are \emph{ad hoc}, relativistic and, contingent. In contrast, Evola envisions a suprantional Imperium, in which the superior power furnishes the structure to the others. This, he claims, is a real unity, opposed to the merely nominal unity defined by Keller.

The final objection is devastating. Keller sees a pacifistic union of separate peoples, each according to its nature. But Evola asks the question, ``what if the nature of a people is to be a warrior race?'' This, Evola claims, is an essential part of Indo-European civilizations. This will blow up the ideal of peaceful ethnocentric enclaves, since one (or more) of them will seek to establish an Empire over the others; this is an essential feature and not a historically contingent event that may or may not occur.

We see the Germanic nations to this day adopting most of Keller's conceptions. The sight of 40,000 Norwegians singing together for peace will not convince another people, whoever they may be, who are of a warrior nature, to leave them in that desired peace. Evola's text follows:
\end{quotationx}


Keller believes in immanent laws of the people, given by ``nature'', respected and followed in themselves, not originating from any power from above or however personified: in any case, having nothing to do with a ``beyond''. This is a ``faith'' like any other. In reality, it would be difficult to adduce a single legal system of ancient peoples, the Aryans included, for whom the authority of the laws was not related to a ``divine'' origin from above, and was not considered to have been introduced by legislators themselves ``divine''. But we already know Keller's point of view: instead of understanding the order proper to an earthly and temporal \emph{Imperium} as the reflection of a transcendent order and, from another point of view, the secularization of law and authority as purely spiritual in origin, he tends to see in every idea of a spiritual \emph{Regnum} a type of fantastic projection of the fantastic image of an earthly reign. This stands, more or less, at the intellectual level of Euhemerism\footnote{\url{http://www.etymonline.com/index.php?term=euhemerism}}.

On the other hand, Keller finds it particularly difficult to define the concept of \emph{Volk}. In certain cases, it seems that he values the national socialist conception of it as the criterion and measure, corresponding to the formula \emph{Volksgemeinschaft}. Where he says that the ``community of nations'', the hoped for supranational order, will not be able to be realized until all the peoples have been realized according to the Nazi totalitarian social concept, expressed precisely by the formula mentioned. But on the other hand, he concedes:

\begin{quotex}
It is not important to see how the individual peoples conceive themselves, even if each one leaves the other free to form for themselves a given concept of their own essence and to live in conformity with it.
\end{quotex}


This is the conclusion of the analysis of the various ideas of the people, a conclusion that evidently leads to a full indifferentism and a pale norm of reciprocal tolerance on the international plane. The fact that Keller insufficiently acknowledges race to define the people; that he introduced the rather indeterminate concept of the ``community of destiny'' (this is a greatness that is defined depending on directions and many directions exist in the history of a great people); that, finally, he speaks of a ``people in becoming'', without being able to say exactly what will define the \emph{terminus ad quem}, i.e., the definite form of such a becoming---all that goes to confirm the mythical and weak content of this conception of the \emph{Volk}. The only coherent solution would be to assert dogmatically that nations as truly such do not exist before being ``total nations'' according to the national socialist formula of the \emph{Volksgemeinschaft}.

Nevertheless, to assume this political form as a condition for the membership with full rights in the \emph{orde nationum} is not, evidently, a falling short of the principle of respect of every ``people'', even to the point of judging them ``as lesser peoples'' and putting them almost under guardianship, even under a foreign State, when it has not reached such a form.

There is more. Keller, as we saw, blames the inorganicity and the leveling that would be typical of every Roman or religious conception of the \emph{Regnum}, or \emph{Reich} if one prefers, of the nations. Such an accusation is even turned against his own conception. In the logic of this, in fact, we don't encounter any obstacle to not end up in the utopian vision of an order that takes up all the peoples of the earth. And we don't encounter any obstacle through the fact that this order is rather more abstract and deprived of content of that of the now defunct League of Nations. In it, would be a question of only taking care of the laws of each people, of teaching them, of maintaining peace and the balance of power---in summary, a type of administrative function that does not presuppose any specific vision of the world, no higher point of reference. But, in its turn, the particularistic conception, relativistic, and collectivistic of the Volk, is the cause of that. If the authority of higher values is not recognized, it is evident that among the various ``national'' unities, only extrinsic relations, ``gravitations'' and ``equilibria'' will be able to exist, which do not involve, at its foundation, anything essential. In every way for us, such an \emph{ordo nationum} as universal order is absurd and undesirable: we conceive of the concrete, differentiated, supranational unity, decentered in a well-defined vision of the world and in well-defined values that give the tone and the internal unity to a given ``imperial space''. The superior civilization of a dominating race must furnish the reference point to a succession of political ethnic minors, in order that these, rising up from simply national values, and integrating such values, they find in them the basis to understand and to feel united to each other: to actually unite, and not in the style of an indifferent tolerance of belonging to the same club.

Therefore, these supranational unities, that will put a stop to the period of particularistic, schismatic, anti-European, as well as differentiated nationalism, will be able to even be combative. We see that Keller acknowledges war, but without recognizing in it any specific value, in the same way that he eventually acknowledges the State, in a transitional phase, as an instrument and pedagogue that may help tie ``people'' to realize themselves, finally, in a ``total'' form. The general intonation of his view remains pacifistic. Wherever he uses the word ``power'', \emph{Macht}, Keller thinks only of oppression, tyranny, violence, distortion of the people. He barely remembers Moltke's phrase: ``universal peace is only a dream, and not even a beautiful dream.''

Nevertheless, since every people, according to him, should be respected according to their own inclinations, we can certainly think that there exist peoples of a warrior race and calling---we noted everything that was said, starting from Klemm and D'Eichtal, about the typological distinction of ``active races'' and ``passive races'', ``conquering races'' and slave races'', etc. Keller, it is true, on the other hand, refers to research on savage tribes, where he would demonstrate natural inclinations to reciprocal respect and peace, for the purpose of challenging the idea that the natural state is the war of all against all. But much more documentation would be necessary to show such a thesis, against which, all the great Aryan history already remains: it is true, not a history of eternal war and for war itself, but history, nevertheless, which is consistent with virile and dominating natures, capable of realizing superior values in battle--- superior, often, to all that can come from a climate of peace of naturalistic harmony. Among the rights recognized by Keller for the people, he acknowledges the right to ``their own development'': it would therefore behoove us to ask why this development must be limited to the bourgeois domain of ``culture and economy'', with the exclusion of everything that refers to the ``power'' factor, where the natural inclination of such people was exactly that of the warrior and heroic type. We think instead that the heroic and warrior elements are of particular importance even as the common basis of an ``imperial space'', i.e., of a concrete and well-articulated order between a given group of nations.

\paragraph{Part III} 

\begin{quotex}
The organization of a block of the type like the European, Aryan, and Roman block is the only concrete task of our future and the only object of a serious consideration in regards to a new law and a new ordering of the people.
\flright{\textsc{Julius Evola}}
\end{quotex}

\begin{quotationx}
In this final installment of the essay, please note these points.

\begin{itemize}


\item Evola declares that Tradition is Roman, it forms a union with preceding traditions, with the Hyperborean tradition as its beginning.

\item The Volk is a modern invention, opposed to Tradition. Its roots are found in Rousseau and the French Revolution.

\item The life and traditions of early Germanic tribes is poorly documented, relying mostly on some comments of Tacitus, and the Eddas, which were preserved much later by post-pagans.

\item However, since the Germans were Aryans, we can surmise about them, based on what we know of other Aryan civilizations. The Roman tradition, particularly that of the Middle Ages, is much more fully documented. Hence, that is where we put our focus.

\item The State organizes the people and relates them to spiritual reality. However, this is not a superstructure whose purpose is only to aggrandize its leaders and oppress the masses, as some in the New Right believe.

\item This orientation to transcendent values is a feature not only of the Christian Middle Ages, but also the authentic pagan traditions.

\item Some civilizations may expand, creating a supranational organization. But this is not merely an administrative function, hegemony, or acting as the ``world police''. Rather, it is to spread its vision of a spiritual, transcendental reality. Blocks of nations will arise, united by common transcendental values.
\end{itemize}
\end{quotationx}


Finding it necessary, therefore, to draw a conclusion, only what refers to the polemic against an abstract juridical internationalistic normativism\footnote{\url{http://fr.wikipedia.org/wiki/Normativisme}} is acceptable in Keller's ideas, after all, up to this point that is what the idea of State has assumed in certain cases in the modern world, in full settlement and against depersonalized ``neutral'' forms, deprived of a substrate of quality and true strength. The points of positive reference invoked by Keller against what are nevertheless as problematic as ever, and hardly acceptable from the point of view of a tradition, that we without difficulty declare to be Roman, in order to immediately add that it in its Romanity forms a unity with the traditions characteristic of every great past civilization, with that of ancient Aryanity at the head. In fact, as we said, this mythology of the \emph{Volk}, of the people-nation as absolute source of every authority, is only a modern invention, that arose in Germany recently and that in Germany itself it now seems somewhat passé, through the same force of European events.

Keller's and other's efforts to relate it back to an ancient German tradition are artificial. We saw, as much, that Keller was constrained to condemn as non-German and ``Romanizing'' the very tradition of the Prussian State and to extend an analogy, he blamed the Medieval Empire. Perhaps he will be able to support himself on such a view of the primitive Germans. But such views interest us very little, because they cannot be authoritative: The ``primitive'' Germans are in fact only some isolated, dispersed and involute stocks of the Great Aryan tree, and we have to judge and understand their traditions on the basis of that which, in a rather more complete, precise, and documented form, present to us the great Aryan civilizations of antiquity, including the Roman.

In such ancient traditions, there are so many ``myths'', but of a more elevated and real content than the modern myth of the ``\emph{Volk}'', which, as we noted, was essentially given to traditional and dynastic Europe by French Jacobinism and has a strongly collectivistic and anti-hierarchical flavor. According to our traditional point of view, the ``nationalities'' can exist but on a naturalistic plane, not yet political nor properly spiritual: in place of certain spontaneous forms of sensibility and of certain customs where, however, the primary element ends up being less the ``nation'' than the race. The centre and the indispensable condition of the nation as political and spiritual reality is instead the State. The State is not a mere fact of ``power'' or a type of abstract juridical superstructure, but an ethical and spiritual reality and a formative and discriminatory force. This force, in its turn, is connected both to an elite, to a race of leaders and to the prestige proper to them, both to transcendent, to a certain extent, values (which are not necessarily only those of the Christian religion: ancient Iran, ancient Rome, etc., teach it) that legitimate them, because there is true authority only as authority from above.

The formative power of such principles, after having unified and animated a ``people'', can be projected beyond its borders and, in different ways, does not exclude those of war, can give rise to higher unities, i.e., supranational, but nevertheless well defined, and ordered by a determinate law: they are the ``imperial spaces''. The ``rank'' that goes to assure to a given nation the supranational directive function cannot be measured in materialist, administrative, bourgeois pacifist terms, and as quasi-police for collective security and peace. On the contrary, it is adherence to already noted transcendental values and to the corresponding faculty of animating, transporting, making capable of energy and commitment, capable even of confronting tragedy and misfortune, which is the sign of such a superiority. The logic of the system does not then lead to its extension to every nation of the world, but to its limitations and particular zones, as blocks of nations, blocks united by the same chain of an ``Order'', capable, where it occurs, of the same heroic unanimity of a ``crusade''.

Leaving aside both universalistic and nationalistic myths, \emph{the organization of a block of the type like the European, Aryan, and Roman block is the only concrete task of our future and the only object of a serious consideration in regards to a new law and a new ordering of the people.}



\flrightit{Posted on 2012-04-24 by Aeneas}

\begin{center}* * *\end{center}

\begin{footnotesize}
\begin{sffamily}
\texttt{Cologero on 2012-04-26 at 08:50 said:}

Evola makes clear the difference between volkisch movements an authentic Tradition. It is absurd to make the people, or the folk, the ultimate judge of anything. According to Evola, it is the elite, by which he apparently means the ruling elite, who speak for the masses and animate them.

We can agree with that, as far as it goes, but, following the Master Rene Guenon, we assert that there is a spiritual elite beyond the ruling elite. The latter's task is to cooperate in disseminating the spiritual vision to the mass, through its rule and positive law.

Evola accepts the Catholic pessimistic view of human nature, but only for the mass. He believes that the ruling elite are different, a position he calls ``realism''. We think instead it is naive. Perhaps somewhere, at some time, the elite had achieved the Primordial State, but certainly not in recent memory.

\hfill

\texttt{Dominion on 2012-05-03 at 17:04 said:}

Reading this, I'm reminded of two things:

The first is the French Emperor Napoleon. He used his power to spread the values of Liberalism across Europe and attempt to create a new order among the nations, even a new aristocracy. Apparently there can be an antitraditional ``Imperium'' as well.

The second is the European Federation envisioned by Guillaume Faye, where ``Evola is reconciled with Marinetti.''

But how, I wonder, is it possible to begin to create such an order? Certainly, the neofascist terrorist attacks of the 1970's did nothing to further the Evolian vision of certain members of that movement. Moreover, any hint of decline in the current order leads people to believe that we are not following the principles of the order itself (Liberalism, democracy, etc) properly. There seems to be no real chance in the near future of a spiritual order in the line of this essay being able to rise. How would you see this happening, given the current climate?

\hfill

\texttt{Cologero on 2012-05-03 at 20:57 said:}

Evola claimed that Freemasonry was able to reorganize itself (after the chaos of Revolution) during Napoleon's reign, eventually leading to the formation of a Masonic Republic in France. As we shall see, despite appearances, that was not a source of order. Faye's ``visions'' are of no interest to us, as they are prone to change every few years. There is no reconciliation between truth and error; men are willing to make such distinctions.

How can order arise from disorder? That fantasy we leave to biologists and free market economists. Yet we have the answer to your question which was given by the writers of Tradition. Does anyone take it seriously?

Guenon says it must begin with an intellectual conversion, that is, the knowledge of principles. If an elite of those who have undergone a conversion can congeal, there is a possibility of change. Evola explains that every great Western civilization began with a divine revelation; that is pretty much the same as what Guenon claimed, only that it is publicly acknowledged. That is, the elite through their Power and Will, can shape and form the people.

Either that thought rings true to you, and it motivates you, or it does not. You will never ``see'' this happening, unless you are among those are ``making'' this happen.

\hfill

\texttt{Graham on 2012-05-03 at 23:18 said:}

``That is, the elite through their Power and Will, can shape and form the people.''

To prevent misunderstanding, it should be pointed out that the formal cause here is inspired. The elite aren't just making it up as they go along.

\hfill

\texttt{Cologero on 2012-05-03 at 23:45 said:}

Thanks, Graham, for making that clear. It is an elite that is chosen from above, not those who manage to grab power. That is what distinguishes them from the likes of a Napoleon, as Dominion mentioned.

\hfill

\texttt{Dominion on 2012-05-04 at 02:46 said:}

``How can order arise from disorder?''

Isn't the whole process of Initiation one from evolving from a disordered and finite ego to a True Will and Eternal Self? Perhaps I misunderstand, but that seems to be exactly what is being described. The answer would then be that the path from disorder to order is one of following the correct process and being guided by the proper laws and forces.

\hfill

\texttt{Graham on 2012-05-04 at 11:13 said:}

Dominium, in saying that initiation is guided by ``the proper laws and forces'' you answered your own question, since he couldn't be guided by an order that didn't already exist in principle.

\hfill

\end{sffamily}
\end{footnotesize}

